
\section{Proof of Proposition~1}
\label{app:proof_interaction_gap}

\begin{proof}
Let $y_i=\Omega(\tilde{s}_i)$, $u_i=u^\pi_H(\tilde{s}_i,a;g)$, and $\hat u_i=h_\theta(y_i,a;g)$.
By the triangle inequality, $\Delta = |u_1-u_2| \leq |u_1-\hat u_1| + |\hat u_1-\hat u_2| + |\hat u_2-u_2|$.
Since $h_\theta$ is $L$-Lipschitz in $y$, $|\hat u_1-\hat u_2| \leq L\|y_1-y_2\| \leq L\epsilon$.
Therefore, $|u_1-\hat u_1|+|u_2-\hat u_2| \geq \Delta-L\epsilon$.
The left-hand side is nonnegative, so it is at least $(\Delta-L\epsilon)_+$.
At least one of the two terms must therefore be at least $(\Delta-L\epsilon)_+/2$.
\end{proof}


\section{Additional results}
\label{app:tables}

This appendix holds the detailed tables behind the main-paper figures: per-task success, the exact terminal-outcome distribution, the full dataset-support distances, the backward-embedding geometry probe, and the matched-observation aliasing numbers.

\subsection{Per-task success}

\begin{table}[H]
\centering
\caption{State-based evaluation success by task (IQM $\pm$ half-width of the $95\%$ stratified-bootstrap CI across $10$ seeds).}
\label{tab:pertask}
\begin{tabular}{lcccccc}
\toprule
\textbf{Method} & \textbf{task1} & \textbf{task2} & \textbf{task3} & \textbf{task4} & \textbf{task5} & \textbf{Overall} \\
\midrule
\fb{}    & $16.7 \pm 10.0$ & $68.3 \pm 10.9$ & $75.0 \pm 10.0$ & $40.0 \pm 13.3$ & $43.3 \pm 18.4$ & $50.0 \pm 7.0$ \\
\gciql{} & $73.7 \pm 5.9$  & $70.0 \pm 7.0$  & $73.0 \pm 7.0$  & $58.0 \pm 5.2$  & $64.3 \pm 7.0$  & $68.3 \pm 3.7$ \\
\crl{}   & $24.0 \pm 3.3$  & $28.0 \pm 6.7$  & $28.0 \pm 4.8$  & $24.3 \pm 3.5$  & $26.3 \pm 5.8$  & $25.9 \pm 2.0$ \\
\rldp{}  & $28.3 \pm 10.8$ & $16.7 \pm 6.7$  & $35.0 \pm 15.9$ & $13.3 \pm 12.5$ & $18.3 \pm 10.0$ & $22.7 \pm 6.0$ \\
\bottomrule
\end{tabular}
\end{table}

\subsection{Terminal-outcome distribution}

\begin{table}[H]
\centering
\caption{Cross-method terminal-outcome distribution (\% of episodes, mean across $10$ seeds; the first five columns sum to $100\%$). The post-contact share is $(\text{maintain}+\text{transport})/(100-\text{success})$.}
\label{tab:cross_method_outcomes}
\begin{tabular}{lcccccc}
\toprule
\textbf{Method} & \textbf{Success} & \textbf{Approach} & \textbf{Grasp} & \textbf{Maintain} & \textbf{Transport} & \textbf{Post-contact share} \\
\midrule
\fb{}    & $49.4$ & $0.0$ & $9.8$  & $23.8$ & $17.0$ & $80.6\%$ \\
\gciql{} & $69.2$ & $0.2$ & $12.2$ & $9.2$  & $9.2$  & $59.7\%$ \\
\crl{}   & $23.2$ & $0.4$ & $23.8$ & $28.6$ & $24.0$ & $68.5\%$ \\
\rldp{}  & $25.4$ & $0.2$ & $9.6$  & $44.0$ & $20.8$ & $86.9\%$ \\
\bottomrule
\end{tabular}
\end{table}

\subsection{Dataset-manifold support}

\begin{table}[H]
\centering
\caption{Nearest-neighbor distance $m_{\mathrm{dist}}$ from rollout states to comparable offline states (mean across seeds; larger means lower support), by per-state regime and terminal outcome. S = success, Mf = maintain-grasp failure, Tf = transport failure. Figure~\ref{fig:support} plots the Mf$-$S and Tf$-$S deltas for the contact regimes.}
\label{tab:state_coverage}
\begin{tabular}{llccc}
\toprule
\textbf{Method} & \textbf{Regime} & \textbf{S} & \textbf{Mf} & \textbf{Tf} \\
\midrule
\fb{}    & approach & $2.28$ & $2.17$ & $2.19$ \\
         & pickup   & $2.23$ & $2.74$ & $2.42$ \\
         & carry    & $2.18$ & $2.31$ & $2.14$ \\
\addlinespace
\gciql{} & approach & $2.50$ & $2.76$ & $2.96$ \\
         & pickup   & $2.93$ & $3.75$ & $3.68$ \\
         & carry    & $2.65$ & $3.00$ & $2.99$ \\
\addlinespace
\crl{}   & approach & $2.53$ & $2.87$ & $2.74$ \\
         & pickup   & $3.08$ & $4.50$ & $3.47$ \\
         & carry    & $2.76$ & $2.88$ & $3.07$ \\
\addlinespace
\rldp{}  & approach & $2.37$ & $2.30$ & $2.42$ \\
         & pickup   & $2.46$ & $2.93$ & $2.90$ \\
         & carry    & $2.21$ & $2.62$ & $2.55$ \\
\bottomrule
\end{tabular}
\end{table}

\subsection{Backward-embedding geometry probe}

\begin{table}[H]
\centering
\caption{Carry-regime resolution of the backward embedding $B$ for cube-to-goal distance $d$, under a linear and a nonlinear (MLP) readout, with the maintain-grasp failure rate. Restricted to the methods with an explicit backward embedding. Better geometric resolution in $B$ does not imply better post-contact control.}
\label{tab:geometry_control}
\begin{tabular}{lccc}
\toprule
\textbf{Method} & \textbf{$R^2(B\to d)$ linear} & \textbf{$R^2(B\to d)$ MLP} & \textbf{Maintain-grasp fail (\%)} \\
\midrule
\fb{}   & $0.20$ & $0.21$          & $23.8$ \\
\rldp{} & $0.55$ & $\mathbf{0.75}$ & $44.0$ \\
\bottomrule
\end{tabular}
\end{table}

\subsection{Representation separability}

The aliasing ladder (Figure~\ref{fig:repsep_ladder}) is the per-task linear-probe
AUC for predicting success-bound versus fail-bound among in-hand (grasp/carry)
training states, probed separately per task (fixed goal) and averaged over the
five tasks ($0.5$ is chance). Table~\ref{tab:repsep_raw} gives the raw-state
ceiling and Table~\ref{tab:repsep_learned} the learned representations: the
goal-independent state embedding $\phi(s)$ (for \gciql{}, the penultimate
goal-conditioned value features; for \crl{}, the state-side contrastive
embedding) and, for the methods with a forward map, the on-policy successor
representation $F(s,\pi,z)$.

\begin{table}[H]
\centering
\caption{Raw-state ceiling: per-task linear-probe AUC from raw state sub-vectors (method-independent; the success-bound label is the play-data reach-before-end flag). The instantaneous configuration only weakly determines the outcome, and the gripper alone is at chance.}
\label{tab:repsep_raw}
\begin{tabular}{lcc}
\toprule
\textbf{Raw feature subset} & \textbf{dims} & \textbf{AUC} \\
\midrule
Gripper (opening + contact) & $2$  & $0.50$ \\
Gripper + arm joints        & $14$ & $0.61$ \\
Full state                  & $28$ & $0.64$ \\
\bottomrule
\end{tabular}
\end{table}

\begin{table}[H]
\centering
\caption{Learned in-hand representations: per-task train AUC for success- vs fail-bound. $\phi(s)$ matches but never exceeds the raw $0.64$ ceiling, and the forward map $F(s,\pi,z)$---built to predict successor occupancy---does no better. Eval-time probing gives the same picture but is high-variance over ${\sim}50$ episodes (AUC swings ${\pm}0.2$ across rollout draws), so we report the stable per-task training-state probe. Figures~\ref{fig:repsep} and~\ref{fig:repsep_tsne} show the corresponding scatters.}
\label{tab:repsep_learned}
\begin{tabular}{lcc}
\toprule
\textbf{Method} & \textbf{$\phi(s)$} & \textbf{$F(s,\pi,z)$} \\
\midrule
\fb{}    & $0.64$ & $0.59$ \\
\gciql{} & $0.66$ & --     \\
\crl{}   & $0.73$ & --     \\
\rldp{}  & $0.64$ & $0.61$ \\
\bottomrule
\end{tabular}
\end{table}

\begin{figure}[H]
\centering
\includegraphics[width=\linewidth]{fig_repsep.png}
\caption{Per-method PCA of the learned in-hand representation: training states (top) and held-out rollouts (bottom); green success-bound, orange fail-bound, with class centroids. Train annotations are per-task AUC; the eval row is a single-draw episode-grouped AUC (high-variance, see Table~\ref{tab:repsep_learned}).}
\label{fig:repsep}
\end{figure}

\begin{figure}[H]
\centering
\includegraphics[width=\linewidth]{fig_repsep_tsne.png}
\caption{As Figure~\ref{fig:repsep} but with a per-panel t-SNE projection (nonlinear; robustness check). Success- and fail-bound grasps remain intermixed for the \fb{} family. The AUC, computed on the full representation, is projection-independent.}
\label{fig:repsep_tsne}
\end{figure}

\begin{figure}[H]
\centering
\includegraphics[width=\linewidth]{fig_value_aliasing.png}
\caption{Value-function aliasing. Top: per-task $z$-scored value $Q=F(s,\pi,z)\cdot z$ for success-bound vs fail-bound in-hand training grasps---the distributions coincide (AUC $0.51$/$0.54$). Bottom: value at the policy's own in-hand grasps split by terminal outcome (eval; illustrative, high-variance over ${\sim}50$ episodes). The value the policy follows does not separate the grasps.}
\label{fig:value_aliasing}
\end{figure}

\subsection{Matched-observation aliasing}

\begin{table}[H]
\centering
\caption{Matched-observation aliasing test: for each fail-bound state, the fraction of nearest success-bound states (in observation space) that the value ranks higher, by regime (mean $\pm$ std across seeds). $0.5$ is chance. Matched pairs lie within the data spread (standardized distance ${\approx}2.7$ vs ${\approx}7$ median). Figure~\ref{fig:aliasing} plots these values.}
\label{tab:aliasing}
\begin{tabular}{lccc}
\toprule
\textbf{Method} & \textbf{Approach} & \textbf{Pickup} & \textbf{Carry} \\
\midrule
\fb{}    & $0.59 \pm 0.01$ & $0.57 \pm 0.01$ & $0.59 \pm 0.01$ \\
\gciql{} & $0.50 \pm 0.01$ & $0.48 \pm 0.01$ & $0.50 \pm 0.01$ \\
\crl{}   & $0.55 \pm 0.02$ & $0.55 \pm 0.01$ & $0.55 \pm 0.02$ \\
\rldp{}  & $0.59 \pm 0.01$ & $0.57 \pm 0.01$ & $0.60 \pm 0.01$ \\
\bottomrule
\end{tabular}
\end{table}
